\documentclass[11pt]{article}
\usepackage[utf8]{inputenc}
\usepackage[T1]{fontenc}
\usepackage{booktabs}
\usepackage{amsmath}
\usepackage{geometry}
\geometry{margin=2.5cm}

\title{Report di Causal Fairness Analysis}
\author{Generato da FairMind + LLM}
\date{2026-08-11}

\begin{document}
\maketitle

\section{Configurazione dell'analisi}

\begin{tabular}{ll}
\toprule
Dataset & adult \\
Attributo protetto ($X$) & S2\_gender \\
Confronto & Female $\to$ Male \\
Outcome ($Y$) & T\_income (>50K) \\
Mediatore ($W$) & hours-per-week \\
Confondente ($Z$) & education \\
\bottomrule
\end{tabular}

\section{Effetti causali (FairMind, valori esatti)}

I cinque effetti seguenti sono calcolati da FairMind tramite inferenza esatta
sulla rete bayesiana stimata. Questi valori sono definitivi e non vanno
ricalcolati o alterati.

\begin{tabular}{lr}
\toprule
Effetto & Valore \\
\midrule
Total Variation (TV) & 0.193714 \\
Total Effect (TE) & 0.184736 \\
Spurious Effect (SE) & 0.008977 \\
Direct Effect (DE) & 0.138404 \\
Indirect Effect (IE, forma inversa) & -0.046333 \\
\bottomrule
\end{tabular}

\section{Interpretazione qualitativa}

\subsection{Effetto totale e componente spuria (TV, TE, SE)}

Il totale disparità tra i due gruppi è di 0.193714, di cui 0.184736 è un effetto causale genuino (TE) che sopravvive alla rimozione del confondente Z. La componente spuria (SE) è di 0.008977, che è causata dal confondente Z piuttosto che dall'attributo protetto X stesso.

\subsection{Effetto diretto (DE)}

Il 0.138404 del totale effetto è diretto, non passando attraverso il mediatore. Questo indica una discriminazione diretta contro il gruppo protetto, con un'entità significativa rispetto al totale effetto causale.

\subsection{Effetto indiretto (IE)}

L'effetto indiretto attraverso il mediatore è di -0.046333, che contrasta con l'effetto diretto. Questo suggerisce che il canale del mediatore agisce in modo a mitigare l'effetto diretto.

\section{Recap Questions}

\begin{enumerate}
\item Esiste una discriminazione diretta nei confronti del gruppo protetto? \textbf{Risposta:} SI
\item L'effetto indiretto attraverso il mediatore mitiga l'effetto totale? \textbf{Risposta:} SI
\item L'effetto totale osservato supera la soglia di rilevanza pratica? \textbf{Risposta:} SI
\item La componente spuria (SE) contribuisce in modo sostanziale alla differenza osservata? \textbf{Risposta:} NO
\item L'effetto diretto è la componente dominante rispetto all'effetto indiretto? \textbf{Risposta:} SI
\end{enumerate}

\end{document}